\section{Guarantees of related attention methods}
\label{app:related}

This appendix states the guarantees of three attention methods
summarized in Section~\ref{sec:prior-main}, in the notation of their
papers. No proof in this paper uses them.

\paragraph{KDEformer.}
KDEformer gives a high-probability spectral approximation to an
attention matrix product \cite{kdeformer2023}. In its notation,
with unnormalized attention matrix $A$, diagonal row-sum matrix
$D$, and value matrix $V$, its guarantee is
\[
 \|D^{-1}AV-\widetilde D^{-1}A\Pi^\top\Pi V\|_{\rm op}
 \le \epsilon\|D^{-1}A\|_{\rm op}\|V\|_{\rm op}.
\]
For $N$ rows, its sample count is
$m=O(\epsilon^{-2}\log N\,[\operatorname{srank}(D^{-1}A)
+\operatorname{srank}(V)])$, where
$\operatorname{srank}(M)=\|M\|_F^2/\|M\|_{\rm op}^2$.
Constructing $\Pi$ and $\widetilde D$ requires the KDE work;
materializing the approximate product adds $O(Nmd)$ operations
at width $d$. The theorem concerns all rows.

\paragraph{HyperAttention.}
HyperAttention proves a related operator-norm guarantee under
conditions on the maximum column norm of normalized attention
and on residual row masses after a heavy-entry mask is supplied
\cite{hyperattention2023}. Under its bounds
$N\max_j\|(D^{-1}A)_{*,j}\|_2^2\le N^{o(1)}$ and
residual row-mass ratio at most $N^{o(1)}$, with
$\epsilon>N^{-o(1)}$, the time is
$O(dN^{1+o(1)}+d\,\operatorname{nnz}(\mathrm{mask}))$
and the success probability is at least $0.98$.
Finding the mask requires its own assumptions.
The guarantee supplies no certified interval for every output
probability on every random run.

\paragraph{MagicPIG.}
MagicPIG uses a prefill-built locality-sensitive hash index for
decoding \cite{magicpig2024}. Its ideal attention-sampling theorem
assumes draws from the true attention distribution and gives an
unbiased sample average. The implemented LSH construction instead
uses inclusion probabilities $u_i$ in the ratio
\[
 \widehat v=
 \frac{\sum_{i\in S}e^{a_i}v_i/u_i}
      {\sum_{i\in S}e^{a_i}/u_i}.
\]
For $L$ tables of $K$ hash bits, retrieving a key after at least
two table collisions gives
$u_i=1-(1-\zeta_i^K)^L-L\zeta_i^K(1-\zeta_i^K)^{L-1}$,
where $\zeta_i$ is its one-bit collision probability.
Such ratios are generally biased. For example, take
weights $(1,2)$, values $(0,1)$, and one uniformly included
index. The self-normalized estimate is the selected value, with
mean $1/2$, while the true weighted mean is $2/3$.
This distinction concerns exactness; empirical approximation
quality requires separate measurements.
